光伏出力与负载预测误差、电价波动及储能运行约束共同影响微网购电费用。本文以购电总费用最小为目标，构建日前购电规划与储能控制相结合的优化框架，研究确定性调度、预测误差和日内预报更新下的购电策略，并扩展至波动电价条件。问题二至四的评估期为2025年2月1日至12月31日，共334天。

\textbf{针对问题一}，建立满足供需平衡、储能容量、充放电功率及日初末储电量相等约束的\textbf{线性规划（LP）模型}。证明在电价非负、允许处理富余电量且切换费用为0时，LP存在满足充放电互斥约束的最优解，并用混合整数线性规划（MILP）和网格动态规划（DP）核验。全天购电量为59,482.70\,kWh，最低费用为35,126.95元，较无储能方案\textbf{节省26.90\%}。

\textbf{针对问题二}，仅用零点前的历史数据，采用岭回归（Ridge）作为固定预测方法，并进行近期偏差修正。根据单时段无储能模型在5倍紧急购电价格下的一阶最优条件，取同小时净负载误差0.8分位数的非负部分作为加在预测负载上的\textbf{用电安全余量}，结合LP确定零点购电策略。进一步采用\textbf{边际价值DP}，权衡当前放电收益与保留电量的后续价值。主方案总费用为1383.2797万元；在冻结同一普通购电计划、保持相同初始及年末储电量的对照中，边际价值DP较贪心控制\textbf{节省3.1750万元}。

\textbf{针对问题三}，融合发布的光伏预报与历史预测，建立计及\textbf{双向策略调整}费用的滚动购电模型。采用“\textbf{全日规划、六小时更新}”机制，每次规划至日末，日内仅确认随后六小时的购电调整，并利用\textbf{递减价值DP}控制储能。总费用为1320.7966万元。在其余设置一致、年末储电量相同的对照中，使用日内更新预报较始终沿用零点预报\textbf{节省6.7227万元}，降幅为0.5064\%，紧急购电量由48,425.57\,kWh降至41,443.14\,kWh。

\textbf{针对问题四}，仅利用决策日前的历史电价训练直方图梯度提升模型，电价预测平均绝对误差为0.0399元/kWh。将问题二、三的策略扩展至波动电价条件，\textbf{按预测电价规划、实际交付电价结算}。在各自实际年末储电量下，问题四-2与问题四-3的总费用分别为1457.6456万元和1394.6759万元，两套完整策略的费用\textbf{相差62.9697万元}。

逐区间核验表明，所得策略满足供需平衡、储电量递推及储能运行约束，各项费用复算一致。问题四-3的储电量网格由50\,kWh细化至25\,kWh后，总费用仅变化49.50元。
